\section{本章小结}
\label{sec:integ-summary}

本章开头问：怎样按位置给出不同的贡献再累计，累计能否与极限、依次求积和求导交换。积分按值域分层：非负简单函数的积分是“高度乘以测度”之和，与写法无关；一般非负可测函数的积分是下方简单函数积分的上确界，允许取$+\infty$，并且任何一列递增的简单逼近都能算出它。实值函数把正部与负部分开累计，只要至少一部分有限积分就存在，两部分都有限才称可积；复值函数以模的积分有限为可积条件，按实部与虚部积分。积分看不见零测集，只依赖函数的几乎处处等价类。

基本性质各有条件。非负函数的线性组合要求系数非负，并按$0\cdot(+\infty)=0$处理无穷；可积函数允许任意实系数或复系数，但只知道积分存在不足以相加。单调性可以放宽到几乎处处，复值积分不比较大小，绝对值估计则通过旋转把复积分化为实部与模的比较，并给出$\lvert\int f-\int g\rvert\leq\int\lvert f-g\rvert$。有限个不交区域的积分可加，增删零测集不改变积分，可数个区域的情形要用单调收敛定理。

与经典积分相比，有界闭区间上的有界函数黎曼可积当且仅当不连续点集零测，此时两种积分相等；$\mathbf 1_{\mathbb Q}$只有新积分。反常积分在绝对收敛时与新积分一致，条件收敛的$\sin x/x$却不是Lebesgue可积的，因为测度积分不依赖任何次序。

积分与极限的交换由尖峰例子引出：逐点收敛不足以保证积分收敛，质量可以挤进越来越窄的区间，或者摊到无界区域上。单调收敛定理处理递增的非负函数列，并给出非负级数逐项积分；Fatou引理对非负函数列给出单向估计；支配收敛定理在存在共同可积上界时交换极限与积分，并保证$\int\lvert f_n-f\rvert\to0$。

累次积分与含参积分分别给出换序与参数运算的条件。$\sigma$有限测度的乘积测度由矩形上的乘积赋值唯一确定；Tonelli定理对非负可测函数无条件换序，Fubini定理在先核对$\int\lvert f\rvert<+\infty$后对一般函数换序，截面只需几乎处处可积，例外处的内层积分取零；完备化下截面的可测性也只对几乎每个参数成立。含参积分的连续性与积分号下求导都依赖参数邻域上与参数无关的可积控制；积分区域移动时，还要加上边界移动带来的项。

微积分基本定理在绝对连续条件下联系积分与导数。可积函数的变上限积分绝对连续，几乎处处可导，导数几乎处处等于被积函数；反过来，绝对连续函数几乎处处可导、导数可积，并能由一点的函数值和导数的积分恢复。Cantor函数连续、单调、几乎处处导数为零却从$0$增到$1$，标出了恢复公式的边界。第~\ref{chap:smoothing-weak-derivatives}章在这些基础上研究卷积光滑化与弱导数，那里要用到本章的换序、积分号下求导与一维的积分恢复；换测度与沿映射换元分别在第~\ref{chap:measure-densities}章与第~\ref{chap:pushforward-measures}章展开。
